\section{Agent 的工作方式：Macro Actions 与并行化}

\subsection{Peter Steinberg 模式}

Karpathy 特别提到了 Peter Steinberg 的工作方式：他同时运行多个 Codex agent session，每个 session 执行约 20 分钟的高难度任务。他在多个 repo 之间切换，给每个 agent 分配独立的功能性任务。

\begin{importantbox}{Macro Actions 思维}
关键的思维转换是：不再以"一行代码""一个函数"为操作单位，而是以"一个功能模块""一个研究方向"为宏操作单位。Agent 1 做功能 A，Agent 2 做功能 B（两者不冲突），Agent 3 做 research，Agent 4 写实现方案——一切以宏操作（macro actions）的方式管理代码仓库。
\end{importantbox}

\begin{figure}[H]
\centering
\includegraphics[width=0.82\textwidth]{frame-macro-actions.jpg}
\caption{视频约 00:04:26：Karpathy 解释 code agent 时代的软件仓库管理单位已从微观改动转向 macro actions}
\end{figure}

\subsection{Skill Issue：一切都是技能问题}

当 agent 失败时，Karpathy 倾向于认为这是"skill issue"而非能力不足——是自己没有给出足够好的指令、没有配置好 agents.md 文件、没有搭建合适的 memory 工具。这种"一切皆可优化"的心态既令人兴奋，也令人焦虑。

\begin{warningbox}{Agent 的 Jaggedness（参差不齐）}
Karpathy 同时描述了 agent 令人沮丧的一面：他觉得自己同时在和"一个极其聪明的 PhD + 一个 10 岁小孩"对话。Agent 在 on-rails 的领域（如通过 RL 训练的编程任务）表现出超人能力，但在 off-rails 领域（如理解用户意图、讲笑话）则停滞不前。这种 jaggedness 是当前模型的根本特征。
\end{warningbox}


